% titel: Blatt 4, A1: Photonen (Lampe) und Aussagen zur Elektronenstruktur
% skript: Mortimer Kap. 6 (Licht, Photonen, Elektronenstruktur)

% Teil I: Werkzeuge, kein Lösungsweg.

Zwei Grundtypen einer Klausurfrage (5 min): eine Rechenaufgabe zur Zahl der
Photonen und eine Wissensabfrage zu Aussagen über die Elektronenstruktur.

\begin{werkzeugkasten}[Photonenenergie und Leistung]
Energie eines Photons aus Frequenz bzw.\ Wellenlänge:
\begin{equation}
  E_{\text{photon}} = h\,\nu = \frac{h\,c}{\lambda},
\end{equation}
mit $h = \qty{6.626e-34}{J\,s}$, $c = \qty{2.998e8}{m/s}$.
Strahlungsleistung $P$ (Energie pro Zeit) liefert in der Zeit $t$ die Energie
\begin{equation}
  E = P\,t.
\end{equation}
Zahl der Photonen $N = E / E_{\text{photon}} = P\,t\,\lambda /(h\,c)$.

\textbf{Konzeptfragen:} Emission = Übergang von angeregt in den Grundzustand
unter Energiefreisetzung; Heisenberg: Ort und Impuls nicht gleichzeitig
scharf bestimmbar; im Atom ist die Aufenthaltswahrscheinlichkeit der
Elektronen \emph{nicht} überall gleich; das Modell ``Elektron im Kasten''
erklärt die Energiequantisierung.

\textit{Skript:} Mortimer Kap. 6 (Licht und Atom), Übung 4/1.
\end{werkzeugkasten}

\begin{achtung}[Darauf achten]
\begin{itemize}
  \item Wellenlänge in Meter umrechnen ($\lambda = 450\ \mathrm{nm}
        = \qty{4.50e-7}{m}$), sonst stimmt $E = hc/\lambda$ nicht.
  \item $E = P\,t$ erst bilden (hier $P = 5.0\ \mathrm{J/s}$) und dann durch
        die Photonenenergie teilen --- nicht Molina-Verwechseln von Energie und
        Leistung.
  \item Bei den Aussagen: ``überall gleiche Wahrscheinlichkeit'' ist falsch;
        das ist die typische Falle.
  \item Auf die Einheiten achten: $h$ in J\,s, $c$ in m/s $\Rightarrow$
        $E$ in J.
\end{itemize}
\end{achtung}
